\documentclass[aspectratio=169,10pt]{beamer}
\usetheme[numbering=fraction,progressbar=frametitle]{metropolis}
\usepackage{fontspec}
\setsansfont{Latin Modern Sans}
\usepackage{amsmath,amssymb,mathtools,graphicx,booktabs}
\usepackage{polyglossia}
\setdefaultlanguage{french}
\setbeamertemplate{navigation symbols}{}
\setbeamerfont{frametitle}{size=\large}
\setbeamerfont{normal text}{size=\normalsize}
\setlength{\parskip}{.45em}
\hypersetup{pdfauthor={},pdfsubject={Lecture slides},pdfkeywords={},pdftitle={Fonction inverse}}
\providecommand{\N}{\mathbb{N}}
\providecommand{\Z}{\mathbb{Z}}
\providecommand{\D}{\mathbb{D}}
\providecommand{\Q}{\mathbb{Q}}
\providecommand{\R}{\mathbb{R}}
\providecommand{\C}{\mathbb{C}}
\providecommand{\vect}[1]{\overrightarrow{#1}}
\providecommand{\vi}{\vec{\imath}}
\providecommand{\vj}{\vec{\jmath}}
\providecommand{\vk}{\vec{k}}
\providecommand{\norm}[1]{\left\lVert #1\right\rVert}
\providecommand{\vc}[2]{\begin{pmatrix}#1\\#2\end{pmatrix}}
\providecommand{\deter}[2]{\det\left(#1,#2\right)}
\providecommand{\abs}[1]{\left\lvert #1\right\rvert}
\providecommand{\dd}{\,\mathrm{d}}
\providecommand{\e}{\mathrm{e}}
\providecommand{\Card}{\operatorname{Card}}
\providecommand{\E}{\operatorname{E}}
\providecommand{\V}{\operatorname{V}}
\providecommand{\degre}{^{\circ}}
\providecommand{\un}[1]{\,\mathrm{#1}}
\providecommand{\sci}[2]{#1\times10^{#2}}
\providecommand{\pgcd}{\operatorname{PGCD}}
\providecommand{\ppcm}{\operatorname{PPCM}}
\title{Fonction inverse}
\subtitle{Cours et exemples}
\date{}
\author{}
\begin{document}
\begin{frame}[plain]
\vfill
{\usebeamerfont{title}\inserttitle\par}
\medskip
{\usebeamerfont{subtitle}\insertsubtitle\par}
\medskip
{\usebeamercolor[fg]{progress bar}\rule{\textwidth}{.5pt}}
\vfill
\end{frame}
\begin{frame}{Cours de référence}
\(f'(x)=-\dfrac1{x^2}<0\) pour tout \(x\neq0\). La fonction inverse est strictement décroissante sur chacun des intervalles \(]-\infty,0[\) et \(]0,+\infty[\). Les droites \(x=0\) et \(y=0\) sont des asymptotes.
\end{frame}
\begin{frame}{Cours de référence}
On ne peut pas écrire que la fonction inverse est décroissante « sur \(\R\setminus\{0\}\) » sans préciser les deux intervalles. Par exemple \(-1<1\) mais \(f(-1)=-1<1=f(1)\), ce qui contreditrait la définition globale de la décroissance.
\end{frame}
\begin{frame}{Cours de référence}
Sur \(]0,+\infty[\), \(g(x)=x+\dfrac4x\) vérifie \(g'(x)=1-\dfrac4{x^2}=\dfrac{(x-2)(x+2)}{x^2}\). Le dénominateur est positif ; \(g'\) est négative sur \(]0,2[\) puis positive sur \(]2,+\infty[\). La fonction admet un minimum \(g(2)=4\).
\end{frame}
\begin{frame}{Définition et courbe de la fonction inverse}
\begin{itemize}
\item Déterminer le domaine.
\item Construire la courbe.
\item Utiliser les symétries.

\end{itemize}
\end{frame}
\begin{frame}{Définition et courbe de la fonction inverse}
\begin{columns}[T,onlytextwidth]
\begin{column}{.57\textwidth}
Calculer \(\dfrac1x\) pour \(x=-4,-2,-1,-\dfrac12,\dfrac12,1,2,4\). Placer les points et décrire les symétries. Pourquoi \(x=0\) est-il absent ?
center center
\end{column}
\begin{column}{.4\textwidth}
\centering
\includegraphics[width=\linewidth,height=4.5cm,keepaspectratio]{figures/stmg-04-2.pdf}
\end{column}
\end{columns}
\end{frame}
\begin{frame}{Définition et courbe de la fonction inverse}
La fonction inverse est \(f(x)=\dfrac1x\) sur \(\R\setminus\{0\}\). Sa courbe est une hyperbole symétrique par rapport à l'origine.
\end{frame}
\begin{frame}{Définition et courbe de la fonction inverse}
Reprendre les données de l'activité et rédiger la solution dans l'ordre : définition des objets, propriété utilisée, calcul exact ou approché, puis phrase de réponse. La réponse modèle figure dans l'édition enseignant.
\end{frame}
\begin{frame}{Définition et courbe de la fonction inverse}
\textbf{Conduite en 1 h :} 7 min d'automatismes ; 10 min de mise en commun et trace ; 28 min d'exercices essentiels ; 5 min de correction ; 3 min de billet de sortie.

\textbf{Conduite en 2 h :} ajouter 20 min de recherche de l'activité, 20 min d'institutionnalisation, un second exercice autonome et 10 min de correction comparée.

\textbf{Différenciation :} fournir aux élèves fragiles une fiche à quatre cases (données, modèle, calcul, conclusion). Demander aux élèves rapides de produire un contre-exemple, une justification ou une seconde méthode.
\end{frame}
\begin{frame}{Comportement aux bornes et asymptotes}
\begin{itemize}
\item Décrire le comportement aux bornes.
\item Identifier les asymptotes.
\item Lire une représentation.

\end{itemize}
\end{frame}
\begin{frame}{Comportement aux bornes et asymptotes}
\begin{columns}[T,onlytextwidth]
\begin{column}{.57\textwidth}
À l'aide d'un tableau de valeurs près de \(0\) et pour de grandes valeurs de \(|x|\), décrire \(\dfrac1x\). Traduire chaque comportement par une asymptote et par une phrase.
center center
\end{column}
\begin{column}{.4\textwidth}
\centering
\includegraphics[width=\linewidth,height=4.5cm,keepaspectratio]{figures/stmg-04-3.pdf}
\end{column}
\end{columns}
\end{frame}
\begin{frame}{Comportement aux bornes et asymptotes}
Pour \(f(x)=\dfrac1x\), \(x=0\) est asymptote verticale et \(y=0\) asymptote horizontale. Les branches s'en approchent sans les couper.
\end{frame}
\begin{frame}{Comportement aux bornes et asymptotes}
Reprendre les données de l'activité et rédiger la solution dans l'ordre : définition des objets, propriété utilisée, calcul exact ou approché, puis phrase de réponse. La réponse modèle figure dans l'édition enseignant.
\end{frame}
\begin{frame}{Comportement aux bornes et asymptotes}
\textbf{Conduite en 1 h :} 7 min d'automatismes ; 10 min de mise en commun et trace ; 28 min d'exercices essentiels ; 5 min de correction ; 3 min de billet de sortie.

\textbf{Conduite en 2 h :} ajouter 20 min de recherche de l'activité, 20 min d'institutionnalisation, un second exercice autonome et 10 min de correction comparée.

\textbf{Différenciation :} fournir aux élèves fragiles une fiche à quatre cases (données, modèle, calcul, conclusion). Demander aux élèves rapides de produire un contre-exemple, une justification ou une seconde méthode.
\end{frame}
\begin{frame}{Dérivée et sens de variation}
\begin{itemize}
\item Dériver la fonction inverse.
\item Justifier les variations.
\item Dresser un tableau.

\end{itemize}
\end{frame}
\begin{frame}{Dérivée et sens de variation}
À partir du taux d'accroissement de \(f(x)=\dfrac1x\), retrouver \(f'(x)=-\dfrac1{x^2}\). Étudier son signe et vérifier la cohérence avec le graphique.
\end{frame}
\begin{frame}{Dérivée et sens de variation}
Pour \(x\neq0\), \(f'(x)=-\dfrac1{x^2}<0\). La fonction inverse décroît sur chacun des intervalles \(]-\infty,0[\) et \(]0,+\infty[\).
\end{frame}
\begin{frame}{Dérivée et sens de variation}
Reprendre les données de l'activité et rédiger la solution dans l'ordre : définition des objets, propriété utilisée, calcul exact ou approché, puis phrase de réponse. La réponse modèle figure dans l'édition enseignant.
\end{frame}
\begin{frame}{Dérivée et sens de variation}
\textbf{Conduite en 1 h :} 7 min d'automatismes ; 10 min de mise en commun et trace ; 28 min d'exercices essentiels ; 5 min de correction ; 3 min de billet de sortie.

\textbf{Conduite en 2 h :} ajouter 20 min de recherche de l'activité, 20 min d'institutionnalisation, un second exercice autonome et 10 min de correction comparée.

\textbf{Différenciation :} fournir aux élèves fragiles une fiche à quatre cases (données, modèle, calcul, conclusion). Demander aux élèves rapides de produire un contre-exemple, une justification ou une seconde méthode.
\end{frame}
\begin{frame}{Combinaisons avec des fonctions polynomiales}
\begin{itemize}
\item Dériver une combinaison.
\item Réduire une expression.
\item Organiser une étude.

\end{itemize}
\end{frame}
\begin{frame}{Combinaisons avec des fonctions polynomiales}
Une entreprise produit \(x>0\) centaines d'objets et son coût moyen est \(C(x)=2x+30+\dfrac{200}{x}\). Calculer \(C'(x)\), factoriser son numérateur et chercher le niveau de production minimisant le coût moyen.
\end{frame}
\begin{frame}{Combinaisons avec des fonctions polynomiales}
Pour étudier \(P(x)+\dfrac{k}{x}\), déterminer le domaine, calculer la dérivée, réduire au même dénominateur, étudier son signe puis dresser les variations.
\end{frame}
\begin{frame}{Combinaisons avec des fonctions polynomiales}
Reprendre les données de l'activité et rédiger la solution dans l'ordre : définition des objets, propriété utilisée, calcul exact ou approché, puis phrase de réponse. La réponse modèle figure dans l'édition enseignant.
\end{frame}
\begin{frame}{Combinaisons avec des fonctions polynomiales}
\textbf{Conduite en 1 h :} 7 min d'automatismes ; 10 min de mise en commun et trace ; 28 min d'exercices essentiels ; 5 min de correction ; 3 min de billet de sortie.

\textbf{Conduite en 2 h :} ajouter 20 min de recherche de l'activité, 20 min d'institutionnalisation, un second exercice autonome et 10 min de correction comparée.

\textbf{Différenciation :} fournir aux élèves fragiles une fiche à quatre cases (données, modèle, calcul, conclusion). Demander aux élèves rapides de produire un contre-exemple, une justification ou une seconde méthode.
\end{frame}
\begin{frame}{Signe et factorisation}
\begin{itemize}
\item Factoriser.
\item Construire un tableau de signes.
\item Respecter les valeurs interdites.

\end{itemize}
\end{frame}
\begin{frame}{Signe et factorisation}
Étudier le signe de \(\dfrac{(x-2)(x+3)}x\). Construire le tableau avec les trois valeurs critiques, distinguer zéro et valeur interdite, puis résoudre l'inéquation \(\dfrac{(x-2)(x+3)}x\geq0\).
\end{frame}
\begin{frame}{Signe et factorisation}
Dans un quotient, le signe dépend du numérateur et du dénominateur. Les valeurs interdites sont exclues du tableau, même si un facteur se simplifie formellement.
\end{frame}
\begin{frame}{Signe et factorisation}
Reprendre les données de l'activité et rédiger la solution dans l'ordre : définition des objets, propriété utilisée, calcul exact ou approché, puis phrase de réponse. La réponse modèle figure dans l'édition enseignant.
\end{frame}
\begin{frame}{Signe et factorisation}
\textbf{Conduite en 1 h :} 7 min d'automatismes ; 10 min de mise en commun et trace ; 28 min d'exercices essentiels ; 5 min de correction ; 3 min de billet de sortie.

\textbf{Conduite en 2 h :} ajouter 20 min de recherche de l'activité, 20 min d'institutionnalisation, un second exercice autonome et 10 min de correction comparée.

\textbf{Différenciation :} fournir aux élèves fragiles une fiche à quatre cases (données, modèle, calcul, conclusion). Demander aux élèves rapides de produire un contre-exemple, une justification ou une seconde méthode.
\end{frame}
\begin{frame}{Dérivées de polynômes de degré au plus 3}
\begin{itemize}
\item Dériver un polynôme.
\item Étudier une dérivée.
\item Déduire des variations.

\end{itemize}
\end{frame}
\begin{frame}{Dérivées de polynômes de degré au plus 3}
Pour \(f(x)=2x^3-3x^2-12x+5\), calculer \(f'(x)\), factoriser, étudier son signe et en déduire les variations de \(f\).
\end{frame}
\begin{frame}{Dérivées de polynômes de degré au plus 3}
\((ax^3+bx^2+cx+d)'=3ax^2+2bx+c\). Le signe de la dérivée donne le sens de variation de la fonction.
\end{frame}
\begin{frame}{Dérivées de polynômes de degré au plus 3}
Reprendre les données de l'activité et rédiger la solution dans l'ordre : définition des objets, propriété utilisée, calcul exact ou approché, puis phrase de réponse. La réponse modèle figure dans l'édition enseignant.
\end{frame}
\begin{frame}{Dérivées de polynômes de degré au plus 3}
\textbf{Conduite en 1 h :} 7 min d'automatismes ; 10 min de mise en commun et trace ; 28 min d'exercices essentiels ; 5 min de correction ; 3 min de billet de sortie.

\textbf{Conduite en 2 h :} ajouter 20 min de recherche de l'activité, 20 min d'institutionnalisation, un second exercice autonome et 10 min de correction comparée.

\textbf{Différenciation :} fournir aux élèves fragiles une fiche à quatre cases (données, modèle, calcul, conclusion). Demander aux élèves rapides de produire un contre-exemple, une justification ou une seconde méthode.
\end{frame}
\begin{frame}{Étude complète et représentation}
\begin{itemize}
\item Mener une étude complète.
\item Tracer une courbe cohérente.
\item Interpréter un extremum.

\end{itemize}
\end{frame}
\begin{frame}{Étude complète et représentation}
\begin{columns}[T,onlytextwidth]
\begin{column}{.57\textwidth}
Étudier complètement \(g(x)=x+\dfrac9x\) sur \(]0,+\infty[\) : domaine, dérivée, signe, variations, minimum, comportement aux bornes et esquisse.
center center
\end{column}
\begin{column}{.4\textwidth}
\centering
\includegraphics[width=\linewidth,height=4.5cm,keepaspectratio]{figures/stmg-04-4.pdf}
\end{column}
\end{columns}
\end{frame}
\begin{frame}{Étude complète et représentation}
Une étude complète comporte : domaine, dérivée, signe, variations, limites ou comportement aux bornes, points remarquables et graphique cohérent.
\end{frame}
\begin{frame}{Étude complète et représentation}
Reprendre les données de l'activité et rédiger la solution dans l'ordre : définition des objets, propriété utilisée, calcul exact ou approché, puis phrase de réponse. La réponse modèle figure dans l'édition enseignant.
\end{frame}
\begin{frame}{Étude complète et représentation}
\textbf{Conduite en 1 h :} 7 min d'automatismes ; 10 min de mise en commun et trace ; 28 min d'exercices essentiels ; 5 min de correction ; 3 min de billet de sortie.

\textbf{Conduite en 2 h :} ajouter 20 min de recherche de l'activité, 20 min d'institutionnalisation, un second exercice autonome et 10 min de correction comparée.

\textbf{Différenciation :} fournir aux élèves fragiles une fiche à quatre cases (données, modèle, calcul, conclusion). Demander aux élèves rapides de produire un contre-exemple, une justification ou une seconde méthode.
\end{frame}
\begin{frame}{Bilan et évaluation}
\begin{itemize}
\item Mobiliser domaine, dérivée et signe.
\item Rédiger une étude.
\item Corriger méthodiquement.

\end{itemize}
\end{frame}
\begin{frame}{Bilan et évaluation}
Corriger une étude volontairement erronée de \(h(x)=x+\dfrac4x\) : domaine annoncé \(\R\), dérivée \(1+\dfrac4{x^2}\), tableau sans rupture en \(0\). Réécrire une solution complète.
\end{frame}
\begin{frame}{Bilan et évaluation}
Contrôler séparément les quatre objets : domaine, expression de la dérivée, signe de la dérivée, variations. Une erreur à une étape ne doit pas masquer les suivantes.
\end{frame}
\begin{frame}{Bilan et évaluation}
Reprendre les données de l'activité et rédiger la solution dans l'ordre : définition des objets, propriété utilisée, calcul exact ou approché, puis phrase de réponse. La réponse modèle figure dans l'édition enseignant.
\end{frame}
\begin{frame}{Bilan et évaluation}
\textbf{Conduite en 1 h :} 7 min d'automatismes ; 10 min de mise en commun et trace ; 28 min d'exercices essentiels ; 5 min de correction ; 3 min de billet de sortie.

\textbf{Conduite en 2 h :} ajouter 20 min de recherche de l'activité, 20 min d'institutionnalisation, un second exercice autonome et 10 min de correction comparée.

\textbf{Différenciation :} fournir aux élèves fragiles une fiche à quatre cases (données, modèle, calcul, conclusion). Demander aux élèves rapides de produire un contre-exemple, une justification ou une seconde méthode.
\end{frame}
\end{document}
